\chapter{Evaluation}\label{chapter:evaluation}

This chapter evaluates the monitoring infrastructure through a series of five
experiments of increasing complexity, from a single-host correctness baseline
to a distributed deployment on three physical machines. The experiments are
designed around six research questions:

\begin{itemize}
  \item \textbf{RQ1} Does the monitor detect property violations correctly,
        and with what intrinsic latency and resource overhead?
  \item \textbf{RQ2} Does distributing evaluation over MQTT preserve verdict
        correctness, and what does the transport hop cost?
  \item \textbf{RQ3} Can the framework observe all three ROS~2 interaction
        primitives (topics, services, and actions)?
  \item \textbf{RQ4} Can an unmodified, production-grade navigation stack
        (Nav2) be supervised, and at what overhead?
  \item \textbf{RQ5} Can properties spanning multiple robot traces be
        evaluated as global predicates?
  \item \textbf{RQ6} Does the complete pipeline sustain distributed operation
        across heterogeneous physical machines?
\end{itemize}

Table~\ref{tab:eval-overview} maps the experiments to these questions. Each
experiment reuses the artifacts of the previous one, so the added complexity
is isolated and measurable.

\begin{table}[htb]
  \centering
  \caption{Evaluation experiments and the research questions they address.}
  \label{tab:eval-overview}
  \begin{tabular}{llll}
    \toprule
    \# & Experiment & Complexity added & RQs \\
    \midrule
    E1 & Local correctness baseline & --- (reference point) & RQ1 \\
    E2 & Distributed evaluation over MQTT & transport split; services, actions & RQ2, RQ3 \\
    E3 & Supervising Nav2 & real navigation stack, dynamic obstacles & RQ3, RQ4 \\
    E4 & Fleet-level verification & multiple robots, global predicate & RQ4, RQ5 \\
    E5 & Distributed heterogeneous deployment & three physical machines, LAN & RQ2, RQ6 \\
    \bottomrule
  \end{tabular}
\end{table}

\section{Experimental Setup and Metrics}\label{sec:eval-setup}

All experiments are executed without the graphical dashboard described in
Chapter~\ref{chapter:implementation}. Every run is driven by command-line
commands and predefined YAML configuration files: the collection tier is
started as \texttt{python3 monitor/monitor\_node.py -c <yaml>}, the
evaluation tier as \texttt{python3 monitor/node\_runner.py -c <yaml>}, and
distributed experiments use a plain Mosquitto broker started from a checked-in
configuration file. Simulation front-ends (Gazebo, RViz) are used for
illustration only and are never part of the experiment control flow. The
experiment definitions, launcher scripts, and analysis tooling are stored in
the repository under \texttt{eval/}, one directory per experiment; each run
writes its resolved configuration, raw record and verdict streams, resource
samples, and computed metrics into a timestamped results directory, so every
reported number is reproducible from the archived artifacts.

\paragraph{Metrics.} Four metric groups recur throughout the chapter.

\begin{description}
  \item[Detection latency.] Wall-clock time from the collector receiving the
    triggering message (the \texttt{timestamp} field of the data record) to
    the verdict reaching its exporter. The built-in verdict exporters stamp
    verdicts with the timestamp of the triggering \emph{input} record, so an
    evaluation-owned exporter plugin
    (\texttt{TimestampedVerdictFileExporter}) additionally records the
    emission wall-clock time per verdict line. Verdicts are joined to their
    triggering records through the \texttt{input\_record\_ids} field. On a
    single host both timestamps come from the same clock; for distributed
    runs the residual clock offset after NTP synchronization is reported
    alongside.
  \item[Throughput and loss.] Records per second sustained per source, and
    message loss measured as gaps in the per-source record sequence numbers
    embedded in record identifiers.
  \item[Verdict correctness.] The online verdict stream is compared one-to-one
    against an offline oracle: the recorded input stream is re-evaluated
    against the same property definition, yielding the expected sequence of
    state transitions. Matched, missed, and spurious verdicts give precision
    and recall. Because verdicts carry their triggering record identifiers,
    a match requires not only the correct verdict polarity but also the
    correct triggering record.
  \item[Resource overhead.] CPU utilization and resident set size of the
    monitoring processes, sampled from \texttt{/proc} at 1\,Hz. For the Nav2
    experiments the overhead on the system under test is additionally
    measured by A/B comparison with the monitor disabled.
\end{description}

\section{E1: Correctness and Intrinsic Overhead Baseline}
\label{sec:eval-e1}

The first experiment establishes that the monitor detects property violations
correctly in the simplest possible setting and measures the intrinsic cost of
the pipeline itself, free of any transport: it is the reference point that
the distributed experiments are compared against.

\paragraph{Setup.} A single Ubuntu machine (Ubuntu 24.04, 16 cores, Python
3.12, ROS~2 Kilted) runs two processes: a deterministic stimulus node and one
monitor process. The stimulus publishes \texttt{/cmd\_vel}
(\texttt{geometry\_msgs/msg/Twist}) at 5\,Hz, alternating
\texttt{linear.x} between $1.85$\,m/s and $0.2$\,m/s every three seconds ---
a square wave of known violation and recovery phases. The monitored property
is a commanded speed limit, $\texttt{linear.x} \le 0.5$\,m/s, evaluated
in-process inside the monitor node by a converter--verdict chain
(\texttt{CmdVelSpeedConverter} followed by an edge-triggered
\texttt{ThresholdVerdict}): a verdict is emitted only when the property
changes state, negative on the first violating record and positive on the
first recovering record. All data records and verdicts are written to JSONL
files for the offline oracle.

\paragraph{Results.} Table~\ref{tab:eval-e1} summarizes the reference run of
120\,s; a shorter smoke run produced consistent values. The stimulus produced
40 property state transitions; the monitor emitted exactly 40 verdicts, each
attributable to the first record crossing the threshold, giving precision and
recall of $1.0$. No record sequence gaps occurred. Detection latency from
collector receipt to verdict emission was $0.36$\,ms on average
($p_{95} = 0.46$\,ms, maximum $0.63$\,ms). Steady-state CPU load of the
monitor process was around $1.2$\,\% of one core with a resident set below
74\,MB; only the first sampling interval, covering interpreter start-up and
DDS discovery, exceeded 5\,\%.

\begin{table}[htb]
  \centering
  \caption{E1 reference run (120\,s, single host, no transport).}
  \label{tab:eval-e1}
  \begin{tabular}{lr}
    \toprule
    Metric & Value \\
    \midrule
    Records collected (\texttt{/cmd\_vel}) & 593 (5.01\,Hz sustained) \\
    Record sequence gaps & 0 \\
    Oracle transitions / online verdicts & 40 / 40 \\
    Precision / recall & 1.0 / 1.0 \\
    Detection latency, mean / median & 0.36\,ms / 0.36\,ms \\
    Detection latency, $p_{95}$ / max & 0.46\,ms / 0.63\,ms \\
    Monitor CPU, steady-state mean & 1.2\,\% \\
    Monitor RSS, maximum & 73.5\,MB \\
    \bottomrule
  \end{tabular}
\end{table}

\paragraph{Conclusion.} With respect to RQ1, the monitor is exact at these
rates on a single host: no messages were lost and the verdict stream matched
the oracle without missed or spurious verdicts. The intrinsic pipeline cost
--- sub-millisecond detection latency and a small constant footprint ---
is negligible compared to the network transports introduced in the following
experiments, so latency differences observed there can be attributed to the
transport and deployment structure rather than to the pipeline itself.

% ---------------------------------------------------------------------------
% The following sections are completed incrementally as experiments E2-E5
% are executed; see eval/PLAN.md for the full experiment designs.
% ---------------------------------------------------------------------------

\section{E2: Distributed Evaluation over MQTT}
\label{sec:eval-e2}

The second experiment splits collection and evaluation into separate
processes connected only through an MQTT broker, and widens the observation
scope to all three ROS~2 interaction primitives. Its central claim is
\emph{distribution transparency}: relocating the evaluation stage behind a
message transport must not change the verdicts.

\paragraph{Setup.} On the same machine as E1, the monitor node observes two
topics (\texttt{/cmd\_vel}, \texttt{/odom}), one service
(\texttt{/reset\_pose}, collected through service introspection events), and
one action (\texttt{/fibonacci}, feedback and status phases). All data
records are exported both to a local file and over MQTT to a dedicated
Mosquitto broker, from which an \texttt{rclpy}-free \texttt{node\_runner}
process consumes them. Three deterministic stimulus nodes drive the system:
the E1 speed square wave, a reset-pose robot whose odd service calls
actually reset the odometry while even calls do not, and a Fibonacci action
fixture that sends a goal of order~6 every five seconds.

Correctness is established by three independent checks. First, the speed
property from E1 is evaluated \emph{twice} from the same collected records
--- in-process inside the monitor node and remotely in the runner behind the
MQTT hop --- and the two verdict streams are compared field-by-field on
property identifier, verdict polarity, and triggering-record identifiers.
Second, the remote verdict stream is verified against the offline oracle as
in E1. Third, the service and action observables are checked against their
constructed ground truth: reset-effect verdicts must alternate between
satisfied (odd calls) and violated (even calls), each attributed to one
service-response record and one odometry record, and every action goal must
yield exactly five feedback records and one \textsc{succeeded} status entry.

\paragraph{Results.} In the 120\,s reference run the monitor collected 1454
records across the four sources without sequence gaps. The local and remote
verdict streams for the speed property were identical: 40 verdicts each,
agreeing on property, polarity, and triggering record for every entry. The
remote stream also matched the offline oracle exactly (precision and recall
$1.0$). The reset-effect property produced 39 verdicts alternating between
satisfied and violated, starting with the effective first call, each
correctly attributed to its service-response and odometry records --- a
property that spans two different interaction primitives. The action fixture
sent 23 goals; the monitor recorded $23 \times 5$ feedback records and 23
\textsc{succeeded} status entries, exactly matching the ground truth.
As a secondary observation, detection latency across the localhost MQTT hop
was $2.4$\,ms on average ($p_{95} = 3.1$\,ms), compared to $0.36$\,ms
in-process in E1.

During bring-up, the correctness harness itself exposed a fixture defect:
the speed stimulus also published an odometry stream fixed at the origin,
which made ineffective resets appear effective. The alternation check failed
immediately and the verdict details (\texttt{distance\_to\_origin} $= 0$ on
even calls) identified the cause; after remapping the extra stream away, all
checks passed. This incident illustrates the practical value of
record-attributed verdicts for diagnosing a monitored system.

\paragraph{Conclusion.} With respect to RQ2, distributing evaluation over
MQTT preserved the verdict stream bit-for-bit under the compared fields; the
transport contributes milliseconds of latency but no semantic change. With
respect to RQ3, topics, services, and actions were monitored concurrently in
one session, and service records participated in a cross-primitive property
correlating service responses with subsequent odometry. This local
transparency result is the invariant that the distributed experiment E5
re-tests across physical machines.

\section{E3: Case Study --- Supervising Nav2}
\label{sec:eval-e3}

The third experiment replaces the synthetic stimuli with a production-grade
navigation stack: an unmodified Nav2 system driving a TurtleBot3 through a
custom corridor world in Gazebo while obstacles are spawned into its path at
runtime. The experiment tests whether the monitoring infrastructure can
supervise a realistic robot software stack without any application
modification, and whether its verdicts remain correct when the monitored
behavior is no longer scripted.

\paragraph{Setup.} Gazebo simulates the corridor world; Nav2 (ROS~2 Kilted,
stock \texttt{nav2\_bringup} launch with only the AMCL initial pose preset
in the parameter file) navigates. The monitor observes \texttt{/cmd\_vel}
(\texttt{TwistStamped}), \texttt{/odom} (rate-throttled to 5\,Hz with field
extraction), and the \texttt{/navigate\_to\_pose} action (feedback and
status), exporting over MQTT to the evaluation runner as in E2. A
deterministic mission script drives three goals through the corridor and
logs every event with wall-clock timestamps: goal~A with a clear corridor;
goal~B in the opposite direction while three boxes are spawned into the
robot's path mid-drive; and goal~C after a wall of boxes has blocked the
direct route. Two properties are evaluated: P1, a commanded speed bound
(threshold below cruise speed, so ordinary driving produces
violation/recovery cycles that exercise the oracle); and P2, a per-goal
deadline of 45\,s on the navigation action, evaluated by a new case-study
converter that raises a violation \emph{in flight} on the first action
record past the deadline, or at terminal status otherwise.

Correctness is again established against ground truth that does not depend
on how the robot happens to move: P1 verdicts are checked against the
offline oracle over the recorded records, and P2 verdicts against the
mission log's independently measured goal durations and outcomes.

\paragraph{Results.} In the reference run the monitor collected 14\,409
records (2\,145 velocity commands at 17.8\,Hz, 11\,606 action records at
94.6\,Hz, and throttled odometry) with zero sequence gaps on the unthrottled
sources. The speed property produced ten verdicts, all matching the offline
oracle (precision and recall $1.0$). The goal-deadline property matched the
mission log for all three goals: goal~A (39.7\,s) satisfied; goal~B, slowed
by the spawned obstacles to 62.0\,s, was flagged as a violation exactly at
the 45\,s deadline --- \emph{16.9 seconds before Nav2 itself reported the
goal finished} --- and goal~C satisfied, because Nav2 simply planned around
the blocking wall and reached the goal in 14.4\,s. The odometry stream shows
3\,274 dropped sequence numbers, which here measure the configured
\texttt{RateThrottler} reduction (about 30\,Hz to 5\,Hz) rather than
transport loss --- the monitor-side transformers reduced the exported
odometry volume by roughly a factor of six.

Two observations are worth recording. First, goal~C illustrates why the
correctness methodology deliberately avoids assuming the robot's behavior:
the blocked corridor was expected to cause a failure, but Nav2 found a
detour, and the ground-truth check accepts either outcome as long as the
verdict agrees with what actually happened. Second, the action feedback
stream dominates record volume (Nav2 publishes NavigateToPose feedback at
its controller rate), which motivates throttling it in the multi-robot
experiments.

\paragraph{Conclusion.} With respect to RQ4, an unmodified Nav2 stack was
supervised end-to-end by configuration alone --- two topics and one action
named in YAML, with the same converter/verdict machinery as the synthetic
experiments. With respect to RQ3, the action observable carried a real
safety property rather than mere record collection, and the in-flight
deadline violation demonstrates the practical value of runtime verification
over post-hoc analysis: the verdict was available 17 seconds before the
navigation stack itself reported the outcome.

\section{E4: Fleet-Level Runtime Verification}
\label{sec:eval-e4}

\todo{E4 pending}
This experiment evaluates per-robot properties and a fleet-level minimum
separation predicate across two Nav2-driven robots on crossing routes.
Results will be added when the experiment has been executed.

\section{E5: Distributed Deployment on Heterogeneous Hardware}
\label{sec:eval-e5}

\todo{E5 pending}
The final experiment runs the E4 scenario with the robots, the monitor tier,
and the verdict tier on three physical machines (an Ubuntu workstation, a
Raspberry Pi, and a macOS laptop) communicating over the local network via
MQTT. Results will be added when the experiment has been executed.
